% ---
% id: receiver/lem:fock-functor
% kind: lem
% title: "Functoriality of the Fock lift"
% summary: >-
%   The Fock lift respects composition and adjoints, preserves every
%   photon-number sector, and takes unitaries, contractions and projectors to
%   unitaries, contractions and projectors. The lift of a unitary implements
%   the mode transformation on the annihilation operators.
% prominence: standard
% status: implemented
% uses: [def:fock]
% assumes: []
% refs: []
% provenance: newly-derived
% implements: certiqs_key.receiver.fock:fock_block
% lean:
% ---

The later derivations work block by block on the lift of
Definition~\ref{def:fock}; this lemma collects the properties of the lift they
use.

\begin{lemma}\label{lem:fock-functor}
Let $\mathcal{K}_0$ be finite-dimensional and let $X,Y$ be operators on
$\mathcal{K}_0$. Then, with $\Fock$ as in \eqref{eq:fockfunctor}:
\begin{enumerate}[label=(\roman*),leftmargin=3.0em]
\item $X^{\otimes n}$ preserves $\Symn{n}(\mathcal{K}_0)$, so $\Fock(X)$ is
  well defined, and it is block diagonal with respect to the photon-number
  decomposition: $\Fock(X)\Pnum{n}=\Pnum{n}\Fock(X)$ for every $n$;
\item $\Fock(XY)=\Fock(X)\Fock(Y)$, $\Fock(\id)=\id$, and
  $\Fock(X)^\dagger=\Fock(X^\dagger)$;
\item if $U$ is unitary then $\Fock(U)$ is unitary; if
  $X^\dagger X\preceq\id$ then $\lVert\Fock(X)\rVert\le1$, so $\Fock(X)$ is a
  bounded operator on $\mathcal{F}(\mathcal{K}_0)$;
\item if $P$ is an orthogonal projector then $\Fock(P)$ is the orthogonal
  projector onto $\mathcal{F}(\operatorname{ran}P)\subseteq
  \mathcal{F}(\mathcal{K}_0)$ --- equivalently, with the occupation vectors
  formed in an orthonormal basis of $\mathcal{K}_0$ adapted to
  $\operatorname{ran}P\oplus\ker P$, onto the closed span of those carrying
  no photon outside $\operatorname{ran}P$;
\item if $U$ is unitary then $\Fock(U)$ implements the mode transformation on
  the annihilation operators,
  $\Fock(U)^\dagger a_i\Fock(U)=\sum_jU_{ij}a_j$.
\end{enumerate}
\end{lemma}

\begin{proof}
(i) For a permutation $\pi$ of $n$ factors let $\Pi_\pi$ be the corresponding
unitary on $\mathcal{K}_0^{\otimes n}$. Then $\Pi_\pi X^{\otimes n}=
X^{\otimes n}\Pi_\pi$, because both sides send
$v_1\otimes\cdots\otimes v_n$ to the tensor product of the $Xv_i$ permuted by
$\pi$. Hence $X^{\otimes n}$ commutes with the orthogonal projector
$P_{\Symn{n}}=\frac{1}{n!}\sum_\pi\Pi_\pi$ and therefore preserves its range.
Block diagonality in $n$ is the definition of $\Fock$ as a direct sum.

(ii) On $\Symn{n}$, $(XY)^{\otimes n}=X^{\otimes n}Y^{\otimes n}$ and both
factors preserve $\Symn{n}$, so the restriction of the product is the product
of the restrictions; $\id^{\otimes n}=\id$; and
$(X^{\otimes n})^\dagger=(X^\dagger)^{\otimes n}$, whose restriction to the
invariant subspace $\Symn{n}$ is the adjoint of the restriction. Summing over
$n$ gives the three identities.

(iii) With $U$ unitary, (ii) gives
$\Fock(U)^\dagger\Fock(U)=\Fock(U^\dagger U)=\id$ and likewise on the other
side. With $X^\dagger X\preceq\id$, $\lVert X\rVert\le1$, so
$\lVert X^{\otimes n}\rVert=\lVert X\rVert^n\le1$; restriction to a subspace
does not increase the norm, so every block of $\Fock(X)$ has norm $\le1$ and
the direct sum is bounded by $1$.

(iv) $P^{\otimes n}$ is an orthogonal projector on $\mathcal{K}_0^{\otimes n}$
with range $(\operatorname{ran}P)^{\otimes n}$, and it commutes with
$P_{\Symn{n}}$ by the computation in (i). The product of two commuting
orthogonal projectors is the orthogonal projector onto the intersection of
their ranges, so the restriction of $P^{\otimes n}$ to $\Symn{n}$ is the
orthogonal projector onto
$\Symn{n}(\mathcal{K}_0)\cap(\operatorname{ran}P)^{\otimes n}
=\Symn{n}(\operatorname{ran}P)$. Choosing an orthonormal basis of
$\mathcal{K}_0$ adapted to $\operatorname{ran}P\oplus\ker P$ identifies that
subspace with the span of the occupation vectors supported on
$\operatorname{ran}P$.

(v) Fix the orthonormal basis $\{\ket{e_j}\}$ used in
Definition~\ref{def:fock}. The creation operator satisfies
$a_i^\dagger\ket{\psi}=\sqrt{n+1}\,P_{\Symn{n+1}}(\ket{e_i}\otimes\ket{\psi})$
for $\ket{\psi}\in\Symn{n}$, which is the standard identification of the mode
description with the symmetric-tensor description. For unitary $U$,
$\Fock(U)P_{\Symn{n+1}}=P_{\Symn{n+1}}U^{\otimes(n+1)}$ on
$\mathcal{K}_0^{\otimes(n+1)}$ by (i), hence
\begin{align}
  \Fock(U)\,a_i^\dagger\ket{\psi}
  &=\sqrt{n+1}\;P_{\Symn{n+1}}\big(U\ket{e_i}\otimes U^{\otimes n}\ket{\psi}\big)
  \nonumber\\
  &=\sum_j U_{ji}\,\sqrt{n+1}\;P_{\Symn{n+1}}
      \big(\ket{e_j}\otimes\Fock(U)\ket{\psi}\big)
  =\Big(\sum_jU_{ji}a_j^\dagger\Big)\Fock(U)\ket{\psi},
\end{align}
using $U\ket{e_i}=\sum_jU_{ji}\ket{e_j}$. Therefore
$\Fock(U)a_i^\dagger\Fock(U)^\dagger=\sum_jU_{ji}a_j^\dagger$; taking adjoints
gives $\Fock(U)a_i\Fock(U)^\dagger=\sum_j\overline{U_{ji}}\,a_j$, and
replacing $U$ by $U^\dagger$ gives
$\Fock(U)^\dagger a_i\Fock(U)=\sum_j\overline{(U^\dagger)_{ji}}a_j
=\sum_jU_{ij}a_j$.
\end{proof}
